\chapter{Introduction générale}
\label{chap:introduction}

Cartographier le fond d'un plan d'eau, suivre un panache de pollution ou inspecter un ouvrage immergé sont des tâches que l'on sait confier à un robot sous-marin autonome. Les confier à une flotte de robots peu coûteux plutôt qu'à un engin unique promet une couverture plus rapide et une tolérance aux pannes, mais se heurte à un obstacle physique déterminant : sous l'eau, la radio ne fonctionne pas. L'eau de mer absorbe les ondes électromagnétiques en quelques dizaines de centimètres, ce qui interdit le lien Wi-Fi ou le réseau maillé classique sur lequel repose la robotique en essaim aérienne ou terrestre. Les alternatives acoustiques, seules à porter réellement en immersion, restent coûteuses, lentes et gourmandes en énergie.

Cette contrainte structure l'ensemble du projet Caiman. Elle impose que chaque robot soit capable de conduire seul sa portion de mission, de mémoriser ce qu'il observe, et de n'échanger avec ses pairs que lors des remontées en surface, où une radio bon marché redevient utilisable. Elle impose aussi une électronique embarquée frugale. Un microcontrôleur unique doit y gérer la propulsion, les capteurs de navigation et de bathymétrie, la gestion d'énergie et la liaison radio. La carte doit rester assez compacte pour tenir dans un tube de faible diamètre, et assez simple pour être fabriquée en série à faible coût.

Le problème posé à ce stage découle de cette double exigence. Une carte électronique Caiman existait déjà, mais n'avait jamais été fabriquée ni éprouvée. Il s'agissait donc de reprendre cette conception, d'en corriger les erreurs de schéma et de routage, d'en préparer un dossier de production industriellement recevable, puis de développer les premières briques logicielles embarquées. La question sous-jacente était de savoir jusqu'où l'on peut valider une architecture de flotte communicante avant de disposer du matériel définitif.

C'est à cette question que répondent le simulateur de mission et le banc \emph{hardware-in-the-loop} développés ensuite : ils permettent d'éprouver le protocole de communication, son codage compact et sa sécurisation sur des processeurs réels, sans attendre la réception de la carte finale.

La démarche suivie dans ce rapport est fondée sur la traçabilité des preuves. Une fonctionnalité présente dans le code n'est pas assimilée à une validation physique, et une géométrie acceptée par un processus de fabrication ne suffit pas à démontrer son fonctionnement électrique. Cette distinction est indispensable ici : la carte a été commandée, mais les délais de fabrication ne permettront pas son bring-up avant la soutenance du \DefenseDate.

Le rapport présente d'abord l'environnement du stage et le système Caiman. Il détaille ensuite l'architecture électronique, le travail de correction et la préparation à la fabrication. Les chapitres suivants abordent le logiciel embarqué, la simulation du réseau et le prototype HIL, avant de synthétiser les validations obtenues, les limites et les perspectives.
